\chapter{Controller Design}
\label{ch:controller_design}

\section{Hierarchical Cascade Control Architecture}
The control of a quadrotor represents a classic underactuated nonlinear control challenge. As derived in Chapter \ref{ch:dynamics}, the system has six degrees of freedom—three translational coordinates $(x, y, z)$ and three rotational coordinates represented by the Euler angles $(\phi, \theta, \psi)$—but is governed by only four independent control inputs $U = [u_1, u_2, u_3, u_4]^T$ representing the total vertical thrust and the three body torques, respectively \cite{ref_77}. 

To stabilize and navigate the quadrotor, we adopt a hierarchical, cascaded control topology consisting of an outer translational loop (position control) and an inner rotational loop (attitude control) \cite{ref_78}. The physical justification for this separation relies on the time-scale separation principle: the rotational dynamics are much faster than the translational dynamics. Therefore, the inner loop can be designed to track high-frequency attitude commands generated by the slower outer position loop, which treats the orientation of the vehicle as virtual control inputs to steer the thrust vector \cite{ref_78}.

The dynamic coupling between translation and rotation is established by defining the virtual control inputs $u_x$ and $u_y$ along the horizontal translational axes:
\begin{align}
u_x &= \cos\phi \sin\theta \cos\psi + \sin\phi \sin\psi \\
u_y &= \cos\phi \sin\theta \sin\psi - \sin\phi \cos\psi
\end{align}
Using these virtual inputs, the global translational accelerations from Chapter \ref{ch:dynamics} are simplified to:
\begin{align}
\ddot{x} &= \frac{u_1}{m} u_x \\
\ddot{y} &= \frac{u_1}{m} u_y \\
\ddot{z} &= \frac{u_1}{m} (\cos\phi \cos\theta) - g
\end{align}

\subsection{Attitude Command Inversion Mapping}
The outer loop position controller calculates the virtual inputs $u_x$ and $u_y$ along with the total thrust $u_1$ to track the spatial coordinates $(x_d, y_d, z_d)$. To achieve translation, these virtual commands must be mapped to desired attitude references—$\phi_d$ (desired roll) and $\theta_d$ (desired pitch)—which are then tracked by the inner loop controller. Assuming a desired yaw heading of $\psi_d$ is specified, we perform an analytical inversion of the rotation matrix to extract the desired angles:
\begin{align}
\phi_d &= \arcsin(u_x \sin\psi_d - u_y \cos\psi_d) \\
\theta_d &= \arcsin\left( \frac{u_x \cos\psi_d + u_y \sin\psi_d}{\cos\phi_d} \right)
\end{align}
This formulation ensures that the lateral tracking forces are physically translated into tilt commands.

\subsection{Second-Order Command Filters}
Because the outer loop controller updates at a lower frequency and can generate step changes in the desired attitude angles $(\phi_d, \theta_d)$, taking direct numerical derivatives of these commands is highly susceptible to sensor noise and communication latency \cite{ref_74}. To generate smooth desired angular velocities ($\dot{\phi}_d, \dot{\theta}_d$) and angular accelerations ($\ddot{\phi}_d, \ddot{\theta}_d$), we pass the commanded angles through a second-order command filter of the form:
\begin{equation}
\ddot{\eta}_{df} + 2\zeta_f \omega_f \dot{\eta}_{df} + \omega_f^2 \eta_{df} = \omega_f^2 \eta_d
\end{equation}
where $\eta_d \in \{\phi_d, \theta_d\}$, $\eta_{df}$ represents the filtered reference angle, $\zeta_f$ is the damping factor, and $\omega_f$ is the natural frequency of the filter. This filter acts as a low-pass differentiator, providing smooth analytical derivatives for the feedforward terms in our advanced controllers.

\section{Derivation of the Baseline Cascaded PID Controller}
To establish a linear baseline, we first implement a classical cascaded Proportional-Integral-Derivative (PID) controller. For each of the six controlled variables, we define the tracking error as the difference between the desired state and the actual state:
\begin{equation}
e_i(t) = i_d(t) - i(t), \quad \dot{e}_i(t) = \dot{i}_d(t) - \dot{i}(t) \quad \text{for } i \in \{x, y, z, \phi, \theta, \psi\}
\end{equation}

\subsection{Outer Loop Position Controllers}
The altitude control input $u_1$ is synthesized by summing a gravitational feedforward term with a three-term PID correction to drive the vertical error $e_z$ to zero:
\begin{equation}
u_{1} = \frac{m}{\cos\phi \cos\theta} \left( \ddot{z}_d + g + K_{p,z} e_z + K_{i,z} \int_{0}^{t} e_z(\tau) d\tau + K_{d,z} \dot{e}_z \right)
\end{equation}
For the horizontal axes $x$ and $y$, the virtual inputs $u_x$ and $u_y$ are designed to establish the required lateral thrust projections:
\begin{align}
u_x &= \frac{m}{u_1} \left( \ddot{x}_d + K_{p,x} e_x + K_{i,x} \int_{0}^{t} e_x(\tau) d\tau + K_{d,x} \dot{e}_x \right) \\
u_y &= \frac{m}{u_1} \left( \ddot{y}_d + K_{p,y} e_y + K_{i,y} \int_{0}^{t} e_y(\tau) d\tau + K_{d,y} \dot{e}_y \right)
\end{align}

\subsection{Inner Loop Attitude Controllers}
The torque control inputs $u_2$ (roll torque), $u_3$ (pitch torque), and $u_4$ (yaw torque) are synthesized using the diagonal moments of inertia $J_x, J_y, J_z$ as scaling factors:
\begin{align}
u_2 = \tau_{\phi} &= J_x \left( \ddot{\phi}_d + K_{p,\phi} e_{\phi} + K_{i,\phi} \int_{0}^{t} e_{\phi}(\tau) d\tau + K_{d,\phi} \dot{e}_{\phi} \right) \\
u_3 = \tau_{\theta} &= J_y \left( \ddot{\theta}_d + K_{p,\theta} e_{\theta} + K_{i,\theta} \int_{0}^{t} e_{\theta}(\tau) d\tau + K_{d,\theta} \dot{e}_{\theta} \right) \\
u_4 = \tau_{\psi} &= J_z \left( \ddot{\psi}_d + K_{p,\psi} e_{\psi} + K_{i,\psi} \int_{0}^{t} e_{\psi}(\tau) d\tau + K_{d,\psi} \dot{e}_{\psi} \right)
\end{align}
While this linear decoupled controller is effective for stable hover operations, it cannot actively reject parametric variations or unmodeled aerodynamic disturbances during aggressive maneuvers.

\section{Conventional Sliding Mode Control (SMC) Derivation}
To introduce robustness against matched perturbations, we derive a first-order Conventional Sliding Mode Controller (SMC) \cite{ref_77}. For each degree of freedom, we define a linear sliding surface $s_i(t)$ as a linear combination of the tracking error and its derivative:
\begin{equation}
s_i(t) = \dot{e}_i(t) + c_i e_i(t)
\end{equation}
where $c_i > 0$ is the sliding surface gain. To prove closed-loop stability, we define a positive-definite Lyapunov candidate function:
\begin{equation}
V_i(t) = \frac{1}{2} s_i(t)^2
\end{equation}
Differentiating this function along the system's trajectories yields:
\begin{equation}
\dot{V}_i(t) = s_i(t) \dot{s}_i(t) = s_i(t) (\ddot{e}_i(t) + c_i \dot{e}_i(t))
\end{equation}
To ensure the sliding surface is attractive, we apply the reachability condition $\dot{V}_i(t) \leq -\eta_i |s_i(t)|$, where $\eta_i > 0$. The control laws are developed by combining the nominal equivalent control $u_{\text{eq}}$ (obtained by setting $\dot{s}_i = 0$ under nominal, disturbance-free conditions) with a discontinuous switching term $u_{\text{sw}} = -W_i \text{sgn}(s_i(t))$.

\subsection{Altitude Subsystem SMC Law}
Evaluating the altitude error dynamics $e_z = z - z_d$:
\begin{equation}
\dot{s}_z = \ddot{z} - \ddot{z}_d + c_z \dot{e}_z = \frac{u_1}{m}\cos\phi\cos\theta - g - \ddot{z}_d + c_z \dot{e}_z
\end{equation}
Setting $\dot{s}_z = -W_z \text{sgn}(s_z) - k_z s_z$ (where $k_z > 0$ is a proportional reaching gain) and solving for the total thrust $u_1$ yields the conventional SMC law:
\begin{equation}
u_1 = \frac{m}{\cos\phi \cos\theta} \left( g + \ddot{z}_d - c_z \dot{e}_z - W_z \text{sgn}(s_z) - k_z s_z \right)
\end{equation}

\subsection{Horizontal Translation SMC Laws}
Applying the same formulation to the virtual inputs $u_x$ and $u_y$:
\begin{align}
u_x &= \frac{m}{u_1} \left( \ddot{x}_d - c_x \dot{e}_x - W_x \text{sgn}(s_x) - k_x s_x \right) \\
u_y &= \frac{m}{u_1} \left( \ddot{y}_d - c_y \dot{e}_y - W_y \text{sgn}(s_y) - k_y s_y \right)
\end{align}

\subsection{Rotational Subsystem SMC Laws}
Evaluating the attitude tracking dynamics and compensating for the rotational Coriolis terms:
\begin{align}
u_2 = \tau_{\phi} &= J_x \left( \ddot{\phi}_d - c_{\phi} \dot{e}_{\phi} - \frac{J_y - J_z}{J_x}\dot{\theta}\dot{\psi} - W_{\phi}\text{sgn}(s_{\phi}) - k_{\phi}s_{\phi} \right) \\
u_3 = \tau_{\theta} &= J_y \left( \ddot{\theta}_d - c_{\theta} \dot{e}_{\theta} - \frac{J_z - J_x}{J_y}\dot{\phi}\dot{\psi} - W_{\theta}\text{sgn}(s_{\theta}) - k_{\theta}s_{\theta} \right) \\
u_4 = \tau_{\psi} &= J_z \left( \ddot{\psi}_d - c_{\psi} \dot{e}_{\psi} - \frac{J_x - J_y}{J_z}\dot{\phi}\dot{\theta} - W_{\psi}\text{sgn}(s_{\psi}) - k_{\psi}s_{\psi} \right)
\end{align}

While these conventional SMC laws provide robust trajectory tracking in the presence of matched disturbances, the discontinuous signum function causes high-frequency chattering in the control commands, which degrades actuator performance and excites unmodeled high-frequency dynamics.

\section{Super-Twisting Sliding Mode Control (STSMC) Derivation}
To mitigate the chattering associated with conventional SMC while maintaining its robust properties, we implement Levant's Super-Twisting Algorithm (STA) \cite{ref_74, ref_76}. Under the STSMC framework, we retain the linear sliding surfaces defined in the previous section:
\begin{equation}
s_i(t) = \dot{e}_i(t) + c_i e_i(t)
\end{equation}
However, the reaching dynamics are driven by a continuous second-order sliding mode reaching law that integrates the discontinuous signum function, effectively shifting the high-frequency switching behind an integrator:
\begin{align}
\dot{s}_i(t) &= -k_{1,i} |s_i(t)|^{1/2} \text{sgn}(s_i(t)) + v_i(t) \\
\dot{v}_i(t) &= -k_{2,i} \text{sgn}(s_i(t))
\end{align}
where $k_{1,i}$ and $k_{2,i}$ are positive controller gains tuned to satisfy finite-time convergence under bounded disturbances.

\subsection{Altitude Subsystem STSMC Law}
Equating the sliding surface derivative $\dot{s}_z$ to the super-twisting reaching law:
\begin{equation}
\frac{u_1}{m}\cos\phi\cos\theta - g - \ddot{z}_d + c_z \dot{e}_z = -k_{1,z} |s_z|^{1/2} \text{sgn}(s_z) + v_z
\end{equation}
Solving for the total thrust command $u_1$ yields:
\begin{equation}
u_1 = \frac{m}{\cos\phi \cos\theta} \left( g + \ddot{z}_d - c_z \dot{e}_z - k_{1,z} |s_z|^{1/2} \text{sgn}(s_z) + v_z \right)
\end{equation}
with the continuous integral term:
\begin{equation}
v_z(t) = \int_{0}^{t} -k_{2,z} \text{sgn}(s_z(\tau)) d\tau
\end{equation}

\subsection{Horizontal Translation STSMC Laws}
\begin{align}
u_x &= \frac{m}{u_1} \left( \ddot{x}_d - c_x \dot{e}_x - k_{1,x} |s_x|^{1/2} \text{sgn}(s_x) + v_x \right), \quad \dot{v}_x = -k_{2,x} \text{sgn}(s_x) \\
u_y &= \frac{m}{u_1} \left( \ddot{y}_d - c_y \dot{e}_y - k_{1,y} |s_y|^{1/2} \text{sgn}(s_y) + v_y \right), \quad \dot{v}_y = -k_{2,y} \text{sgn}(s_y)
\end{align}

\subsection{Rotational Subsystem STSMC Laws}
Evaluating the attitude dynamics under the super-twisting reaching law:
\begin{align}
u_2 = \tau_{\phi} &= J_x \left( \ddot{\phi}_d - c_{\phi} \dot{e}_{\phi} - \frac{J_y - J_z}{J_x}\dot{\theta}\dot{\psi} - k_{1,\phi}|s_{\phi}|^{1/2}\text{sgn}(s_{\phi}) + v_{\phi} \right), \quad \dot{v}_{\phi} = -k_{2,\phi}\text{sgn}(s_{\phi}) \\
u_3 = \tau_{\theta} &= J_y \left( \ddot{\theta}_d - c_{\theta} \dot{e}_{\theta} - \frac{J_z - J_x}{J_y}\dot{\phi}\dot{\psi} - k_{1,\theta}|s_{\theta}|^{1/2}\text{sgn}(s_{\theta}) + v_{\theta} \right), \quad \dot{v}_{\theta} = -k_{2,\theta}\text{sgn}(s_{\theta}) \\
u_4 = \tau_{\psi} &= J_z \left( \ddot{\psi}_d - c_{\psi} \dot{e}_{\psi} - \frac{J_x - J_y}{J_z}\dot{\phi}\dot{\theta} - k_{1,\psi}|s_{\psi}|^{1/2}\text{sgn}(s_{\psi}) + v_{\psi} \right), \quad \dot{v}_{\psi} = -k_{2,\psi}\text{sgn}(s_{\psi})
\end{align}

STSMC significantly reduces control chattering by applying a continuous command to the physical motors. However, the use of linear sliding surfaces only guarantees asymptotic (exponential) convergence of the tracking error to zero as time approaches infinity.

\section{Proposed Non-singular Super-Twisting Terminal SMC (NSTSMC)}
To achieve finite-time convergence of tracking errors to absolute zero without inducing high-frequency chattering or mathematical singularities, we propose a Non-singular Super-Twisting Terminal Sliding Mode Controller (NSTSMC) \cite{ref_75}. This controller represents the primary control contribution of this thesis.

\subsection{The Singularity Problem in Terminal SMC}
In traditional Terminal Sliding Mode Control (TSMC), the sliding manifold is designed as:
\begin{equation}
s_{\text{tsm}} = e + \beta \dot{e}^{p/q}
\end{equation}
where $1 < p/q < 2$. Differentiating this surface yields:
\begin{equation}
\dot{s}_{\text{tsm}} = \dot{e} + \beta \frac{p}{q} \dot{e}^{p/q - 1} \ddot{e}
\end{equation}
Setting $\dot{s}_{\text{tsm}} = 0$ to solve for the equivalent control law requires dividing by the velocity term:
\begin{equation}
u_{\text{eq}} \propto -\frac{q}{\beta p} \dot{e}^{2 - p/q}
\end{equation}
Because $1 < p/q < 2$, the exponent is strictly positive. However, when solving for the switching control term, we obtain:
\begin{equation}
u_{\text{sw}} \propto -\dot{e}^{1 - p/q} \text{sgn}(s_{\text{tsm}})
\end{equation}
Because $p/q > 1$, the exponent $1 - p/q$ is negative (e.g., if $p=5, q=3$, then $1 - p/q = -2/3$). Consequently, if the velocity error $\dot{e}$ is zero while the tracking error $e$ is non-zero, the control input grows infinitely, causing a mathematical singularity that saturates the physical actuators and destabilizes flight \cite{ref_75}.

\subsection{The Proposed Non-singular Terminal Sliding Surface}
To eliminate this singularity, we implement a Non-singular Terminal Sliding Mode (NTSM) manifold for each of the six tracking loops \cite{ref_75}:
\begin{equation}
s_i(t) = e_i(t) + \frac{1}{\beta_i} \text{sgn}(\dot{e}_i(t)) |\dot{e}_i(t)|^{\gamma}
\end{equation}
where $\beta_i > 0$ is the sliding surface gain, and the exponent $\gamma = p/q$ is chosen such that $1 < \gamma < 2$. In our implementation, we select $p = 5$ and $q = 3$, yielding $\gamma = 5/3 \approx 1.67$. Differentiating this non-singular sliding surface with respect to time yields:
\begin{equation}
\dot{s}_i(t) = \dot{e}_i(t) + \frac{\gamma}{\beta_i} |\dot{e}_i(t)|^{\gamma - 1} \ddot{e}_i(t)
\end{equation}

By shifting the fractional exponent to the velocity term, the acceleration term $\ddot{e}_i$ can be isolated. To avoid singularity issues and prevent numerical errors associated with fractional powers of negative numbers in MATLAB, we utilize a custom mathematical formulation that expresses the derivative of the surface using an absolute value and signum convention:
\begin{equation}
\ddot{e}_i(t) = -\frac{\beta_i}{\gamma} |\dot{e}_i(t)|^{2 - \gamma} \text{sgn}(\dot{e}_i(t)) - u_{\text{reach},i}(t)
\end{equation}
Because $\gamma = 5/3$, the exponent $2 - \gamma = 2 - 5/3 = 1/3 > 0$ remains strictly positive, ensuring that the control effort is smooth and singularity-free for all state configurations.

\subsection{Synthesizing the Unified NSTSMC Flight Control Laws}
By combining the non-singular terminal sliding manifold with the continuous Super-Twisting reaching law, we obtain the unified NSTSMC laws for all six axes.

\subsubsection{Altitude Control Law ($u_1$)}
For the altitude error $e_z = z - z_d$, the dynamics are:
\begin{equation}
\ddot{e}_z = \ddot{z} - \ddot{z}_d = \frac{u_1}{m}\cos\phi\cos\theta - g - \ddot{z}_d
\end{equation}
Substituting this into the non-singular sliding surface derivative and solving for $u_1$ yields:
\begin{equation}
u_1 = \frac{m}{\cos\phi \cos\theta} \left( g + \ddot{z}_d - \frac{\beta_z}{\gamma} |\dot{e}_z|^{2 - \gamma} \text{sgn}(\dot{e}_z) - k_{1,z} |s_z|^{1/2} \text{sgn}(s_z) + v_z \right)
\end{equation}
with the super-twisting integrator:
\begin{equation}
\dot{v}_z = -k_{2,z} \text{sgn}(s_z)
\end{equation}

\subsubsection{Horizontal Translation Control Laws ($u_x, u_y$)}
Similarly, the virtual control inputs for horizontal tracking are formulated as:
\begin{align}
u_x &= \frac{m}{u_1} \left( \ddot{x}_d - \frac{\beta_x}{\gamma} |\dot{e}_x|^{2 - \gamma} \text{sgn}(\dot{e}_x) - k_{1,x} |s_x|^{1/2} \text{sgn}(s_x) + v_x \right), \quad \dot{v}_x = -k_{2,x} \text{sgn}(s_x) \\
u_y &= \frac{m}{u_1} \left( \ddot{y}_d - \frac{\beta_y}{\gamma} |\dot{e}_y|^{2 - \gamma} \text{sgn}(\dot{e}_y) - k_{1,y} |s_y|^{1/2} \text{sgn}(s_y) + v_y \right), \quad \dot{v}_y = -k_{2,y} \text{sgn}(s_y)
\end{align}

\subsubsection{Rotational Subsystem Control Laws ($u_2, u_3, u_4$)}
The stabilizing torques about the body axes are designed to reject Coriolis cross-coupling:
\begin{align}
u_2 &= J_x \left( \ddot{\phi}_d - \frac{J_y - J_z}{J_x}\dot{\theta}\dot{\psi} - \frac{\beta_{\phi}}{\gamma} |\dot{e}_{\phi}|^{2 - \gamma} \text{sgn}(\dot{e}_{\phi}) - k_{1,\phi}|s_{\phi}|^{1/2}\text{sgn}(s_{\phi}) + v_{\phi} \right) \\
u_3 &= J_y \left( \ddot{\theta}_d - \frac{J_z - J_x}{J_y}\dot{\phi}\dot{\psi} - \frac{\beta_{\theta}}{\gamma} |\dot{e}_{\theta}|^{2 - \gamma} \text{sgn}(\dot{e}_{\theta}) - k_{1,\theta}|s_{\theta}|^{1/2}\text{sgn}(s_{\theta}) + v_{\theta} \right) \\
u_4 &= J_z \left( \ddot{\psi}_d - \frac{J_x - J_y}{J_z}\dot{\phi}\dot{\theta} - \frac{\beta_{\psi}}{\gamma} |\dot{e}_{\psi}|^{2 - \gamma} \text{sgn}(\dot{e}_{\psi}) - k_{1,\psi}|s_{\psi}|^{1/2}\text{sgn}(s_{\psi}) + v_{\psi} \right)
\end{align}
where the auxiliary variables are integrated as:
\begin{equation}
\dot{v}_{\phi} = -k_{2,\phi} \text{sgn}(s_{\phi}), \quad \dot{v}_{\theta} = -k_{2,\theta} \text{sgn}(s_{\theta}), \quad \dot{v}_{\psi} = -k_{2,\psi} \text{sgn}(s_{\psi})
\end{equation}

\subsection{Rigorous Closed-Loop Lyapunov Stability Proof}
To mathematically prove the tracking stability of the proposed NSTSMC law, we formulate a Lyapunov proof for a general second-order dynamic subsystem of the quadrotor \cite{ref_81}:
\begin{align}
\dot{x}_1(t) &= x_2(t) \\
\dot{x}_2(t) &= f(x,t) + g(x,t)u(t) + d(t)
\end{align}
where $d(t)$ represents matched parametric disturbances bounded by $|d(t)| \leq D$. We define the tracking error as $e = x_1 - x_{1d}$, and the non-singular terminal sliding surface as $s = e + \frac{1}{\beta} \dot{e}^{\gamma}$. 

Taking the time derivative of $s$ along the trajectories yields:
\begin{align}
\dot{s} &= \dot{e} + \frac{\gamma}{\beta} |\dot{e}|^{\gamma - 1} \ddot{e} \\
&= \dot{e} + \frac{\gamma}{\beta} |\dot{e}|^{\gamma - 1} (f(x,t) + g(x,t)u(t) + d(t) - \ddot{x}_{1d})
\end{align}
Substituting the proposed NSTSMC control law:
\begin{equation}
u = \frac{1}{g(x,t)} \left( \ddot{x}_{1d} - f(x,t) - \frac{\beta}{\gamma} |\dot{e}|^{2 - \gamma} \text{sgn}(\dot{e}) - k_1 |s|^{1/2} \text{sgn}(s) + v \right)
\end{equation}
where $\dot{v} = -k_2 \text{sgn}(s)$. This yields the closed-loop reaching dynamics:
\begin{equation}
\dot{s} = \frac{\gamma}{\beta} |\dot{e}|^{\gamma - 1} \left( -k_1 |s|^{1/2} \text{sgn}(s) + v + d(t) \right)
\end{equation}

Let us define a state vector for the sliding surface dynamics: $\zeta = [|s|^{1/2}\text{sgn}(s), v]^T = [\zeta_1, \zeta_2]^T$. The derivative of $\zeta$ can be written in matrix form:
\begin{equation}
\dot{\zeta} = \frac{\gamma}{\beta} |\dot{e}|^{\gamma - 1} \frac{1}{2|\zeta_1|} \left( A \zeta + B d(t) \right)
\end{equation}
where the state and input matrices are:
\begin{equation}
A = \begin{bmatrix} -k_1 & 1 \\ -2k_2 & 0 \end{bmatrix}, \quad B = \begin{bmatrix} 1 \\ 0 \end{bmatrix}
\end{equation}
We define a positive-definite Lyapunov candidate function for the sliding surface dynamics \cite{ref_76}:
\begin{equation}
V(\zeta) = \zeta^T P \zeta
\end{equation}
where $P$ is a symmetric, positive-definite matrix:
\begin{equation}
P = \begin{bmatrix} p_1 & p_2 \\ p_2 & p_3 \end{bmatrix} = \begin{bmatrix} 2k_2 + \frac{1}{2} k_1^2 & -\frac{1}{2} k_1 \\ -\frac{1}{2} k_1 & 1 \end{bmatrix}
\end{equation}
Taking the time derivative of $V(\zeta)$:
\begin{align}
\dot{V}(\zeta) &= \dot{\zeta}^T P \zeta + \zeta^T P \dot{\zeta} \\
&= \frac{\gamma}{\beta} |\dot{e}|^{\gamma - 1} \frac{1}{2|\zeta_1|} \zeta^T \left( A^T P + P A \right) \zeta + \frac{\gamma}{\beta} |\dot{e}|^{\gamma - 1} \frac{1}{|\zeta_1|} d(t) B^T P \zeta
\end{align}
By design, the matrix inequality $A^T P + P A \leq -Q$ is satisfied for a symmetric positive-definite matrix $Q$ as long as the super-twisting gains $k_1$ and $k_2$ are chosen to be strictly positive and configured to overpower the disturbance bound $D$ \cite{ref_76}:
\begin{align}
k_1 &> 0 \\
k_2 &> \frac{1}{2} D + \frac{k_1^2 (4D + 1)}{8(k_1 - 2D)}
\end{align}
This guarantees that:
\begin{equation}
\dot{V}(\zeta) \leq -\frac{\gamma}{\beta} |\dot{e}|^{\gamma - 1} \frac{1}{2|\zeta_1|} \zeta^T Q \zeta \leq 0
\end{equation}
Consequently, as long as the tracking velocity error $\dot{e} \neq 0$, the Lyapunov derivative is strictly negative-definite ($\dot{V} < 0$), ensuring that the sliding variable $s$ and its derivative $\dot{s}$ converge to zero in finite time. Once $s=0$ is reached, the tracking error converges to zero along the terminal trajectory in a finite time interval $T_{\text{reach}}$ given by:
\begin{equation}
T_{\text{reach}} \leq \frac{\gamma \beta}{\gamma - 1} |e(0)|^{\frac{\gamma - 1}{\gamma}}
\end{equation}
This completes the stability proof for the proposed NSTSMC architecture.

\section{Summary of Quadrotor Control Laws}
The physical configurations, sliding surfaces, and command structures for all four flight controllers are summarized in Table \ref{tab:controller_comparison}.

\begin{table}[h!]
\centering
\caption{Comparison of Flight Control Strategies and Mathematical Structures}
\label{tab:controller_comparison}
\begin{tabular}{|l|c|c|l|}
\hline
\textbf{Controller} & \textbf{Sliding Surface $s_i$} & \textbf{Reaching Dynamics $\dot{s}_i$} & \textbf{Primary Advantage} \\
\hline
Cascaded PID & N/A & N/A & Linear, simple tuning \\
Conventional SMC & $\dot{e}_i + c_i e_i$ & $-W_i \text{sgn}(s_i) - k_i s_i$ & Robust matched rejection \\
Super-Twisting SMC & $\dot{e}_i + c_i e_i$ & $-k_{1i}|s_i|^{1/2}\text{sgn}(s_i) + v_i$ & Significant chattering reduction \\
Proposed NSTSMC & $e_i + \frac{1}{\beta_i}|\dot{e}_i|^{\gamma}\text{sgn}(\dot{e}_i)$ & $-k_{1i}|s_i|^{1/2}\text{sgn}(s_i) + v_i$ & Finite-time, chattering-free \\
\hline
\end{tabular}
\end{table}

% derivation of PID \\
% maths of smc \\
% derivation of smc \\
% about STA \\
% derivation of stsmc \\
% sigularity issue intro \\
% derivation of nstsmc \\
% final equations of all 4 controllers \\
